\subsection{導入：なぜ数学的基礎がソフトウェアに必要なのか}
ソフトウェアは、曖昧なままでは実行できない。人間同士の会話では「だいたい」「普通は」「適切に」といった表現が許されることがある。しかし、コード、仕様、テスト、モデルでは、どの条件でどの処理を行うかを明確にする必要がある。数学的基礎は、この明確化を支える。

数学は、数だけを扱う学問ではない。より広く見れば、数学は形式体系を扱う。形式体系とは、対象、記号、規則、推論方法を定め、その中で何が成り立つかを調べる枠組みである。ソフトウェア技術者は、業務、通信、画面操作、状態遷移、データ、セキュリティ規則などを形式化しなければならない。

\begin{notebox}{重要}
第17章の数学は、計算式を速く解くためだけの数学ではない。ソフトウェアで「何が正しいか」「どの入力でどの出力になるか」「どの状態からどの状態へ進むか」「どの文字列が有効か」を曖昧なく扱うための数学である。
\end{notebox}
